% ============================================================
%  Multi-Horizon Kirchhoff Residuals: Self-Referential
%  Portfolio Optimality and Transaction-Time Measurement
%  in Hierarchical ETF Systems
% ============================================================
\documentclass[11pt,a4paper]{article}

\usepackage[margin=1.25in]{geometry}
\usepackage{amsmath,amssymb,amsthm}
\usepackage{mathtools}
\usepackage{bm}
\usepackage{bbm}
\usepackage{graphicx}
\usepackage{hyperref}
\usepackage[numbers,sort&compress]{natbib}
\usepackage{microtype}
\usepackage{booktabs}
\usepackage{enumitem}
\usepackage{xcolor}
\usepackage{tikz}
\usetikzlibrary{arrows.meta,positioning,decorations.pathreplacing}

\hypersetup{
  colorlinks=true, linkcolor=blue!60!black,
  citecolor=green!50!black, urlcolor=blue!70!black
}

% ---- theorem environments -------------------------------------------
\theoremstyle{plain}
\newtheorem{theorem}{Theorem}[section]
\newtheorem{lemma}[theorem]{Lemma}
\newtheorem{proposition}[theorem]{Proposition}
\newtheorem{corollary}[theorem]{Corollary}
\theoremstyle{definition}
\newtheorem{definition}[theorem]{Definition}
\newtheorem{remark}[theorem]{Remark}
\newtheorem{example}[theorem]{Example}
\theoremstyle{remark}
\newtheorem{notation}[theorem]{Notation}

% ---- macros ---------------------------------------------------------
\newcommand{\R}{\mathbb{R}}
\newcommand{\E}{\mathbb{E}}
\newcommand{\N}{\mathbb{N}}
\newcommand{\Z}{\mathbb{Z}}
\newcommand{\F}{\mathcal{F}}
\newcommand{\Lop}{\mathbf{L}}        % graph Laplacian matrix
\newcommand{\Lpinv}{\mathbf{L}^{\dagger}} % Moore-Penrose pseudoinverse
\newcommand{\proj}{\Pi_{\Delta}}     % simplex projection
\newcommand{\wstar}{w^{*}}
\newcommand{\Deltam}{\Delta_m}
\newcommand{\lam}{\lambda}
\newcommand{\lamtwo}{\lambda_2}
\newcommand{\lammax}{\lambda_{\max}}
\newcommand{\norm}[1]{\left\|#1\right\|}
\newcommand{\abs}[1]{\left|#1\right|}
\newcommand{\inner}[2]{\langle #1,\, #2\rangle}
\newcommand{\one}{\mathbf{1}}
\newcommand{\kappa_{\gamma}}{\kappa_{\gamma}}
\newcommand{\Gain}{\mathcal{G}}
\newcommand{\GainK}{\Gain_{k}}
\DeclareMathOperator{\tr}{tr}
\DeclareMathOperator{\rank}{rank}
\DeclareMathOperator{\sign}{sign}
\DeclareMathOperator*{\argmin}{arg\,min}

% ---- title ----------------------------------------------------------
\title{%
  \large\bfseries
  Multi-Horizon Kirchhoff Residuals:\\
  Self-Referential Portfolio Optimality and\\
  Transaction-Time Measurement in Hierarchical ETF Systems%
}

\author{%
  \normalsize
  \textit{Anonymous submission for review}%
}

\date{\normalsize\today}

% =====================================================================
\begin{document}
\maketitle

% ---- abstract -------------------------------------------------------
\begin{abstract}
We develop a self-contained mathematical theory of portfolio optimality
that requires no external benchmark.  Starting from the graph-Laplacian
contraction framework of Banach~\cite{Banach1922} and the discrete
current-balance law of Kirchhoff~\cite{Kirchhoff1847}, we construct, for
every prediction horizon $\tau>0$, a unique fixed-point portfolio
$\wstar(\tau)\in\Deltam$ whose node-wise allocation is governed by a
discrete Kirchhoff law on the asset correlation graph.  The residual
between the predicted and realised return at $\wstar(\tau)$ defines a
\emph{Kirchhoff gain-loss process} $\{\GainK(t)\}$ that is a
martingale difference sequence in transaction time---a stochastic clock
driven by the accumulated absolute gain rather than calendar intervals
\cite{Clark1973}.

Three principal consequences follow.  First, Kirchhoff residuals of
systems with different graph structures inhabit incommensurable normed
spaces; no single external scalar (Sharpe ratio, alpha, drawdown) can
serve as a universal performance proxy, a theorem we call \emph{portfolio
incommensurability}.  Second, the distance between fixed points at two
calendar times is bounded above by $\norm{\mu(t)-\mu(t')}_2/\lamtwo(\Lop)$,
so high algebraic connectivity damps allocation drift and stabilises the
entire hierarchy.  Third, nesting the gain-loss trigger into a
hierarchical sequence of stopping times produces a multi-scale
\emph{gear network}---each layer fires only when the layer below has
accumulated a threshold imbalance---whose equilibrium is globally
self-consistent without any external rebalancing signal.

We prove existence, uniqueness, measurability, and almost-sure ergodic
convergence of the transaction-time average portfolio, and we
characterise the fixed-point surface $\tau\mapsto\wstar(\tau)$ as a
Lipschitz curve on $\Deltam$.  The resulting architecture constitutes a
\emph{self-generated ETF}: a basket defined entirely by internal
Kirchhoff equilibrium, passive at every layer and globally
self-correcting.
\end{abstract}

\tableofcontents
\newpage

% =====================================================================
\section{Introduction}
\label{sec:intro}

Classical portfolio theory---from Markowitz's mean-variance frontier
\cite{Markowitz1952} to the continuous-time Hamilton-Jacobi-Bellman
approach of Merton \cite{Merton1969}---is benchmarked against an
external criterion: a market index, a risk-free rate, or a reference
investor.  This dependence is not merely computational; it is
epistemological.  The theory implicitly assumes that \emph{good
performance} has a universal meaning that transcends the structure of
any individual portfolio.

The present paper argues, and proves rigorously, that this assumption
fails when portfolio dynamics are governed by a Laplacian contraction on
a correlation graph.  The reason is structural: each such system
generates a Kirchhoff equilibrium in a Hilbert space whose inner product
depends on the graph Laplacian.  Residuals from different systems
therefore lie in different normed spaces and cannot be compared by a
single scalar.  The only operationally valid optimality criterion is
internal consistency---the Kirchhoff condition is satisfied or it is not.

The contribution of this paper is to make this observation precise and
to build the full mathematical apparatus around it:

\begin{enumerate}[label=(\roman*)]
  \item A \emph{horizon surface} $\tau\mapsto\wstar(\tau)$
        (\S\ref{sec:multihorizon}), a Lipschitz curve on $\Deltam$
        parametrised by prediction horizon $\tau$.

  \item A \emph{Kirchhoff gain-loss process} $\GainK(t,\tau)$
        (\S\ref{sec:gainloss}), the monetary expression of the
        Kirchhoff residual, proved to be a martingale difference
        sequence in the transaction-time filtration.

  \item The \emph{incommensurability theorem}
        (\S\ref{sec:incommensurability}): no external performance
        scalar is invariant under changes of graph structure.

  \item A \emph{fixed-point drift bound} (\S\ref{sec:drift}) that
        quantifies how fast $\wstar$ moves as return forecasts evolve.

  \item A \emph{hierarchical gear network} (\S\ref{sec:gear}) of nested
        stopping times, each layer triggered by the accumulated
        gain-loss imbalance of the layer below, proved to be
        well-defined and a.s.~finite.

  \item \emph{Transaction-time ergodic convergence}
        (\S\ref{sec:ergodic}): the time-average of $\wstar$ along the
        transaction clock converges almost surely to the stationary
        mean under mild mixing conditions.
\end{enumerate}

Throughout, $m$ denotes the number of assets, $G=(V,E,\omega)$ a
weighted undirected graph on $m$ vertices, $\Lop\in\R^{m\times m}$ its
Laplacian, and $\lamtwo$ the Fiedler value (algebraic connectivity)
\cite{Fiedler1973}.  The probability simplex is
$\Deltam=\{w\in\R^m:w\ge 0,\,\one^{\top}w=1\}$.

% =====================================================================
\section{Background and Notation}
\label{sec:background}

\subsection{Spectral graph theory}
\label{sec:specgraph}

Let $G=(V,E,\omega)$ be a connected weighted undirected graph with
vertex set $V=\{1,\dots,m\}$, edge set $E\subseteq\binom{V}{2}$, and
positive edge weights $\omega:\,E\to(0,\infty)$.  The \emph{weighted
adjacency matrix} $A\in\R^{m\times m}$ has entries
$A_{ij}=\omega_{ij}$ for $\{i,j\}\in E$ and $A_{ij}=0$ otherwise.
The \emph{degree matrix} $D=\mathrm{diag}(d_1,\dots,d_m)$ with
$d_i=\sum_j A_{ij}$ is the diagonal matrix of weighted degrees.  The
\emph{graph Laplacian} is
\begin{equation}
  \Lop = D - A.
  \label{eq:laplacian}
\end{equation}
$\Lop$ is symmetric positive semidefinite \cite{Chung1997}, with
spectrum $0=\lam_1\le\lam_2\le\cdots\le\lammax$.  Connectivity of $G$
is equivalent to $\lamtwo>0$ \cite{Fiedler1973,Mohar1991}.

The \emph{Moore-Penrose pseudoinverse} $\Lpinv$ \cite{Penrose1955} of
$\Lop$ satisfies $\Lop\Lpinv=\Lpinv\Lop=I-\frac{1}{m}\one\one^{\top}$,
the orthogonal projector onto $\one^{\perp}$, and
$\norm{\Lpinv}_2=1/\lamtwo$.

\subsection{Contraction mappings and Banach's theorem}
\label{sec:banach}

A mapping $T:(X,d)\to(X,d)$ on a complete metric space is a
\emph{$\kappa$-contraction} if $d(Tx,Ty)\le\kappa\,d(x,y)$ for all
$x,y\in X$ and some $\kappa\in[0,1)$.  Banach's fixed-point theorem
\cite{Banach1922} guarantees that $T$ has a unique fixed point
$x^*\in X$ and that the iterates $x^{(n+1)}=Tx^{(n)}$ satisfy
\begin{equation}
  d\bigl(x^{(n)},x^*\bigr) \;\le\; \kappa^n\,d\bigl(x^{(0)},x^*\bigr)
  \;\xrightarrow{n\to\infty}\; 0.
  \label{eq:banach_rate}
\end{equation}

\subsection{Simplex projection}
\label{sec:simplexproj}

The Euclidean projection $\proj:\R^m\to\Deltam$ is non-expansive:
$\norm{\proj v-\proj u}_2\le\norm{v-u}_2$ for all $u,v\in\R^m$
\cite{BauschkeCombettes2011}.  An explicit $O(m\log m)$ algorithm via
the sorted-threshold construction is given in
\cite{BauschkeCombettes2011}.

\subsection{Martingales and stopping times}
\label{sec:martingales}

Let $(\Omega,\F,\Pr)$ be a probability space with filtration
$(\F_t)_{t\ge 0}$.  A sequence $(X_k)_{k\ge 1}$ adapted to a
sub-filtration $(\F_{t_k})$ is a \emph{martingale difference sequence}
if $\E[X_k\mid\F_{t_{k-1}}]=0$ a.s.~for all $k$ \cite{Doob1953,Williams1991}.
A stopping time $\tau:\Omega\to[0,\infty]$ is \emph{a.s.~finite} if
$\Pr(\tau<\infty)=1$.  The optional stopping theorem and Wald's identity
are standard references in \cite{Doob1953,Shiryaev1978}.

% =====================================================================
\section{Multi-Horizon Fixed-Point Portfolios}
\label{sec:multihorizon}

\subsection{The portfolio update map}
\label{sec:updatemap}

Fix a connected weighted graph $G$ with Laplacian $\Lop$.  For a
prediction horizon $\tau>0$, let $\hat\mu(\tau)\in\R^m$ be the vector
of forecast expected returns over $[t,t+\tau]$, and let
$\hat\mu_c(\tau)=\hat\mu(\tau)-\bar\mu(\tau)\one$ be the
mean-centred version, $\bar\mu(\tau)=\frac{1}{m}\one^{\top}\hat\mu(\tau)$.

\begin{definition}[Portfolio update map]
  For step size $\gamma>0$, define $T_\tau:\Deltam\to\Deltam$ by
  \begin{equation}
    T_\tau(w) \;=\; \proj\!\bigl[(I-\gamma\Lop)\,w + \gamma\hat\mu(\tau)\bigr].
    \label{eq:Ttau}
  \end{equation}
\end{definition}

\begin{theorem}[Multi-Horizon Contraction]
  \label{thm:contraction}
  For any $\gamma\in(0,1/\lammax)$, the map $T_\tau$ is a
  $\kappa_\gamma$-contraction on $(\Deltam,\norm{\cdot}_2)$ with
  \begin{equation}
    \kappa_\gamma \;=\; 1 - \gamma\lamtwo \;\in\; [0,1).
    \label{eq:kappa}
  \end{equation}
  In particular, $T_\tau$ has a unique fixed point
  $\wstar(\tau)\in\Deltam$.
\end{theorem}

\begin{proof}
  Let $w,v\in\Deltam$.  Using non-expansiveness of $\proj$ and
  symmetry of $\Lop$:
  \begin{align}
    \norm{T_\tau(w)-T_\tau(v)}_2
    &\;\le\; \norm{(I-\gamma\Lop)(w-v)}_2 \notag\\
    &\;=\;   \norm{(I-\gamma\Lop)^{1/2}(I-\gamma\Lop)^{1/2}(w-v)}_2 \notag\\
    &\;\le\; \norm{I-\gamma\Lop}_2\,\norm{w-v}_2.
    \label{eq:contraction_step}
  \end{align}
  The spectrum of $I-\gamma\Lop$ lies in $[1-\gamma\lammax,1]$.  Since
  $\Lop\ge 0$ and $\gamma<1/\lammax$, all eigenvalues are positive and
  the spectral norm equals $\max(1-\gamma\lamtwo,\,\gamma\lammax-1)$.
  Because $\gamma<1/\lammax$ implies $\gamma\lammax<1$, the dominant
  eigenvalue is $1-\gamma\lamtwo$, giving
  $\norm{I-\gamma\Lop}_2=1-\gamma\lamtwo=\kappa_\gamma<1$.  Banach's
  theorem \cite{Banach1922} then guarantees a unique fixed point.
\end{proof}

\begin{theorem}[Kirchhoff Portfolio Formula]
  \label{thm:kirchhoff}
  The unique fixed point satisfies
  \begin{equation}
    \wstar(\tau) \;=\; \Lpinv\hat\mu_c(\tau) + \tfrac{1}{m}\one,
    \label{eq:wstar}
  \end{equation}
  and is the unique solution in $\Deltam$ to the discrete Kirchhoff
  law
  \begin{equation}
    \Lop\,\wstar(\tau) \;=\; \hat\mu(\tau) - \bar\mu(\tau)\one.
    \label{eq:kirchhoff_law}
  \end{equation}
\end{theorem}

\begin{proof}
  At a fixed point $w^*=T_\tau(w^*)$, the projection is trivially
  the identity on the active face, so
  $(I-\gamma\Lop)w^*+\gamma\hat\mu(\tau)=w^*$, which gives
  $\Lop w^*=\hat\mu(\tau)$.  Decompose $\hat\mu(\tau)=\hat\mu_c(\tau)+\bar\mu(\tau)\one$
  and recall $\Lop\one=0$, so $\Lop w^*=\hat\mu_c(\tau)$.  The general
  solution to $\Lop x=\hat\mu_c(\tau)$ in $\one^{\perp}$ is
  $x=\Lpinv\hat\mu_c(\tau)$ (minimum-norm solution \cite{Penrose1955}).
  Adding the constraint $\one^{\top}w^*=1$ yields the shift
  $\frac{1}{m}\one$, establishing \eqref{eq:wstar}.  Equation
  \eqref{eq:kirchhoff_law} is the same identity written differently.

  To confirm the fixed-point projection step is trivial, note that
  $\one^{\top}\wstar(\tau)=\one^{\top}\Lpinv\hat\mu_c(\tau)+1=0+1=1$
  (since $\hat\mu_c\in\one^{\perp}$ and $\Lpinv\hat\mu_c\in\one^{\perp}$),
  so the simplex constraint $\one^{\top}w=1$ is automatically
  satisfied.  Non-negativity requires $w^*_i\ge 0$; the projection
  handles the boundary correctly and uniqueness in $\Deltam$ follows
  from strict contractivity.
\end{proof}

\begin{remark}[Financial interpretation of \eqref{eq:kirchhoff_law}]
  At the fixed point, $\sum_j\Lop_{ij}w^*_j(\tau)=\hat\mu_i(\tau)-\bar\mu(\tau)$
  for each $i$.  This means the weighted net flow from asset $i$ to its
  neighbours (correlation-weighted) equals the excess return of $i$ over
  the portfolio average.  Assets with above-average predicted returns
  export weight; assets with below-average predicted returns import
  weight.  The fixed point is a global \emph{Kirchhoff balance}: every
  node is in equilibrium with its graph neighbourhood, exactly as in an
  electrical resistor network \cite{Kirchhoff1847}.
\end{remark}

\subsection{The horizon surface}
\label{sec:horizonsurface}

\begin{definition}[Horizon surface]
  The \emph{horizon surface} is the map
  $\Sigma:\,(0,\infty)\to\Deltam$, $\tau\mapsto\wstar(\tau)$.
\end{definition}

\begin{proposition}[Lipschitz regularity of the horizon surface]
  \label{prop:lipschitz}
  Assume the forecast map $\tau\mapsto\hat\mu(\tau)$ is
  $L_\mu$-Lipschitz in $\R^m$.  Then $\Sigma$ is Lipschitz with
  constant $L_\mu/\lamtwo$:
  \begin{equation}
    \norm{\wstar(\tau)-\wstar(\tau')}_2
    \;\le\; \frac{L_\mu}{\lamtwo}\,\abs{\tau-\tau'}.
    \label{eq:horizon_lipschitz}
  \end{equation}
\end{proposition}

\begin{proof}
  From \eqref{eq:wstar},
  $\wstar(\tau)-\wstar(\tau')=\Lpinv[\hat\mu_c(\tau)-\hat\mu_c(\tau')]$.
  Since $\norm{\Lpinv}_2=1/\lamtwo$ and $\hat\mu_c(\tau)-\hat\mu_c(\tau')$
  is the mean-centred difference (a linear operation of norm at most
  $\sqrt{1-1/m}\le 1$), we obtain
  $\norm{\wstar(\tau)-\wstar(\tau')}_2\le
  (1/\lamtwo)\norm{\hat\mu(\tau)-\hat\mu(\tau')}_2\le
  L_\mu|\tau-\tau'|/\lamtwo$.
\end{proof}

\begin{remark}
  Higher algebraic connectivity (denser, more uniform correlation graph)
  compresses the horizon surface toward a single point.  The Fiedler
  value thus measures how sensitive the optimal basket is to changes in
  the investment horizon.
\end{remark}

% =====================================================================
\section{The Kirchhoff Gain-Loss Process}
\label{sec:gainloss}

\subsection{Definition}
\label{sec:gainloss_def}

Let $(R_i(t,\tau))_{i=1}^m$ denote the realised return of asset $i$
over $[t,t+\tau]$.  The \emph{Kirchhoff gain-loss} of asset $i$ at
time $t$ under horizon $\tau$ is
\begin{equation}
  G_i(t,\tau) \;=\; w^*_i(\tau)\,\bigl[R_i(t,t+\tau)-\hat\mu_i(\tau)\bigr].
  \label{eq:gain_loss}
\end{equation}
The \emph{portfolio-level Kirchhoff gain-loss} is the sum
\begin{equation}
  \Gain(t,\tau) \;=\; \sum_{i=1}^m G_i(t,\tau)
  \;=\; \wstar(\tau)^{\top}\!\bigl[R(t,t+\tau)-\hat\mu(\tau)\bigr],
  \label{eq:total_gain}
\end{equation}
the inner product of the fixed-point weights with the return forecast
error.

\begin{remark}[Monetary interpretation]
  $G_i(t,\tau)$ is the profit or loss attributable to node $i$
  arising from the discrepancy between the predicted return and the
  realised return.  It is zero in expectation when $\hat\mu$ is an
  unbiased forecaster, and its magnitude measures how far the graph
  Kirchhoff condition departs from exact balance ex post.
\end{remark}

\subsection{Martingale structure in transaction time}
\label{sec:martingale}

Following Clark \cite{Clark1973}, we time-change calendar time $t$
by the \emph{transaction clock}
\begin{equation}
  \theta(T) \;=\; \inf\!\left\{t\ge 0:\int_0^t\abs{\Gain(s,\tau)}\,ds > T\right\},
  \label{eq:txclock}
\end{equation}
the right-continuous inverse of the accumulated absolute gain-loss
process.  Let $\{t_k\}_{k\ge 1}$ be the successive transaction times
defined recursively by $t_1=\theta(0)$ and $t_{k+1}=\theta(t_k+\epsilon)$
for a fixed resolution $\epsilon>0$.  Define the discrete filtration
$\F_k=\sigma\bigl(R(\cdot,\cdot),\,\hat\mu(\cdot)\text{ up to }t_k\bigr)$.

\begin{theorem}[Martingale Difference Property]
  \label{thm:martingale}
  Suppose the return process $\{R_i(t,t+\tau)\}$ is adapted to
  $(\F_t)$ and that $\hat\mu_i(\tau)=\E[R_i(t,t+\tau)\mid\F_{t-}]$
  is the conditional mean given the strict past.  Then the sequence
  $\{\Gain(t_k,\tau)\}_{k\ge 1}$ is a martingale difference sequence
  with respect to $(\F_{t_k})$:
  \begin{equation}
    \E\!\bigl[\Gain(t_k,\tau)\mid\F_{t_{k-1}}\bigr] \;=\; 0
    \quad\text{a.s.~for all }k\ge 1.
    \label{eq:mds}
  \end{equation}
\end{theorem}

\begin{proof}
  By the tower property of conditional expectation and measurability
  of $\wstar(\tau)$ with respect to $\F_{t_{k-1}}$ (the fixed-point
  weights depend only on the Laplacian and the forecast $\hat\mu(\tau)$
  formed before time $t_k$),
  \begin{align}
    \E\!\bigl[\Gain(t_k,\tau)\mid\F_{t_{k-1}}\bigr]
    &= \E\!\left[\wstar(\tau)^{\top}(R(t_k,t_k+\tau)-\hat\mu(\tau))
       \mid\F_{t_{k-1}}\right] \notag\\
    &= \wstar(\tau)^{\top}
       \E\!\left[R(t_k,t_k+\tau)-\hat\mu(\tau)\mid\F_{t_{k-1}}\right]
       \notag\\
    &= \wstar(\tau)^{\top}(\hat\mu(\tau)-\hat\mu(\tau)) \;=\; 0.
    \label{eq:mds_proof}
  \end{align}
  Here we used $\F_{t_{k-1}}$-measurability of $\wstar(\tau)$ and
  $\hat\mu(\tau)$ in the second equality and unbiasedness
  $\E[R(t_k,t_k+\tau)\mid\F_{t_{k-1}}]=\hat\mu(\tau)$ in the third.
\end{proof}

\begin{corollary}[Zero-drift accumulation]
  \label{cor:zero_drift}
  Under the conditions of Theorem~\ref{thm:martingale}, the partial
  sums $S_K=\sum_{k=1}^K\Gain(t_k,\tau)$ form a martingale.  In
  particular, $\E[S_K]=0$ for all $K$ and the process has no
  systematic drift.
\end{corollary}

\begin{proof}
  Immediate from Doob's martingale theory \cite{Doob1953,Williams1991}
  applied to the martingale difference sequence.
\end{proof}

\begin{remark}[Transaction time as the natural clock]
  Calendar time $t$ is an arbitrary external reference.  The
  accumulated gain $\int_0^t|\Gain(s,\tau)|ds$ is an intrinsic
  measure of how much information the portfolio has processed.
  Time-changing to the transaction clock $\theta$ converts the
  irregular arrival of information into a uniformly-paced discrete
  sequence, following Clark's subordination programme \cite{Clark1973}.
  The martingale property~\eqref{eq:mds} holds in both time systems,
  but is most naturally interpreted in transaction time because each
  increment $t_{k+1}-t_k$ carries exactly $\epsilon$ units of
  accumulated activity.
\end{remark}

% =====================================================================
\section{Portfolio Incommensurability}
\label{sec:incommensurability}

Classical performance measurement compares portfolios via a scalar:
the Sharpe ratio, Jensen's alpha relative to a market factor, or the
Sortino ratio \cite{Markowitz1952,Merton1969}.  We now prove that
such comparisons are ill-defined when the two portfolios arise from
different graph structures.

\subsection{The Kirchhoff inner product}
\label{sec:kirchhoff_inner}

\begin{definition}[Kirchhoff inner product]
  For a connected graph with Laplacian $\Lop$, define the
  \emph{Kirchhoff inner product} on $\one^{\perp}$ by
  \begin{equation}
    \inner{u}{v}_{\Lop} \;=\; u^{\top}\Lpinv v.
    \label{eq:kirchhoff_inner}
  \end{equation}
  The induced norm is $\norm{u}_{\Lop}=\sqrt{u^{\top}\Lpinv u}$.
\end{definition}

The Kirchhoff inner product is the natural inner product for measuring
residuals of the system $\Lop w=f$: the minimum-norm solution has
$\norm{\wstar}_{\Lop}^2=\hat\mu_c^{\top}(\Lpinv)^2\hat\mu_c$, which
is finite and uniquely determined by $\Lop$.

\begin{theorem}[Portfolio Incommensurability]
  \label{thm:incommensurability}
  Let $(G,\Lop,\hat\mu)$ and $(G',\Lop',\hat\mu')$ be two portfolio
  systems with $G\ne G'$ (non-isomorphic weighted graphs).  Then
  \begin{enumerate}[label=(\roman*)]
    \item\label{it:incomm_norms}
      The Kirchhoff norms $\norm{\cdot}_{\Lop}$ and
      $\norm{\cdot}_{\Lop'}$ are not equivalent in general: there
      exists no finite constant $C>0$ such that
      $\norm{u}_{\Lop}\le C\norm{u}_{\Lop'}$ for all
      $u\in\one^{\perp}\cap\R^m$ (or $\R^{m'}$ if $m\ne m'$).

    \item\label{it:incomm_residuals}
      The Kirchhoff residuals $G(t,\tau)$ and $G'(t,\tau)$ are
      projections of the same return process onto different bases;
      no component-wise or aggregate comparison can identify which
      system is closer to its respective Kirchhoff equilibrium
      without knowledge of both $\Lop$ and $\Lop'$.

    \item\label{it:incomm_benchmark}
      Any scalar benchmark $B\in\R$ used to compare the two
      systems must implicitly assume a common reference Laplacian
      $\Lop_B$; changing $\Lop_B$ reverses the ranking with
      probability one for generic $\hat\mu,\hat\mu'$.
  \end{enumerate}
\end{theorem}

\begin{proof}
  \textbf{Part~\ref{it:incomm_norms}.}  The norm
  $\norm{u}_{\Lop}^2=u^{\top}\Lpinv u$ depends on the eigendecomposition
  of $\Lop$.  If $G\ne G'$, there exists a vector $e\in\one^{\perp}$
  such that $\Lpinv e$ and $(\Lop')^{\dagger}e$ point in different
  directions (by the spectral theorem applied to the distinct Laplacian
  spectra).  The ratio $\norm{e}_{\Lop}/\norm{e}_{\Lop'}$ is therefore
  unbounded as $e$ ranges over $\one^{\perp}$, since one can choose $e$
  aligned with a direction that is expensive under $\Lpinv$ but cheap
  under $(\Lop')^{\dagger}$ \cite{HornJohnson1985}.

  \textbf{Part~\ref{it:incomm_residuals}.}  Write
  $G(t,\tau)=\wstar(\tau)^{\top}(R-\hat\mu(\tau))$ and
  $G'(t,\tau)=(\wstar)'(\tau)^{\top}(R-\hat\mu'(\tau))$.  The fixed
  points $\wstar$ and $(\wstar)'$ are elements of $\Deltam$ and
  $\Delta_{m'}$, respectively; if $m\ne m'$ they are not even elements
  of the same space.  If $m=m'$, the weights are spectral functions
  of $\Lop$ and $\Lop'$ through \eqref{eq:wstar}, and a linear
  comparison $G-G'$ mixes residuals measured in different inner
  products, analogous to subtracting temperatures in Kelvin and
  Fahrenheit without unit conversion.

  \textbf{Part~\ref{it:incomm_benchmark}.}  Suppose the benchmark
  ranks $G$ above $G'$ via $\Gain>\Gain'$.  This comparison
  implicitly uses the Euclidean inner product, i.e.~the Laplacian
  $\Lop_B=I$ (identity graph, all assets identical).  For any
  alternative reference $\Lop_B'$ with $\lamtwo(\Lop_B')\ne 1$,
  the induced inner product reweights the residual coordinates and
  the comparison can be reversed.  Formally, for generic
  $\hat\mu,\hat\mu'$ the set of $\Lop_B$ for which the ranking is
  preserved has measure zero in the space of positive semidefinite
  matrices, a direct consequence of the open mapping theorem in
  $\R^{m\times m}$ \cite{GolubVanLoan1996}.
\end{proof}

\begin{remark}[The epistemological consequence]
  Theorem~\ref{thm:incommensurability} is not a statement about
  computation but about measurement.  Just as there is no universal
  unit for measuring the ``correctness'' of a solution to two
  different linear systems, there is no universal scalar for
  measuring the performance of two portfolios with different graph
  structures.  The Kirchhoff residual is the only internally
  consistent criterion: it measures exactly how far the current
  allocation departs from equilibrium \emph{under the system's own
  Laplacian}.
\end{remark}

% =====================================================================
\section{Fixed-Point Drift Bound}
\label{sec:drift}

In practice, return forecasts $\hat\mu(\tau)$ evolve as new data
arrive.  We now bound how fast the optimal allocation moves.

\begin{theorem}[Fixed-Point Drift Bound]
  \label{thm:drift}
  Let $\wstar(t,\tau)$ denote the fixed-point portfolio computed with
  the forecast $\hat\mu(t,\tau)$ formed at calendar time $t$.  For
  any two calendar times $t,t'$,
  \begin{equation}
    \norm{\wstar(t,\tau)-\wstar(t',\tau)}_2
    \;\le\; \frac{\norm{\hat\mu(t,\tau)-\hat\mu(t',\tau)}_2}{\lamtwo(\Lop)}.
    \label{eq:drift_bound}
  \end{equation}
\end{theorem}

\begin{proof}
  By \eqref{eq:wstar},
  \begin{align}
    \wstar(t,\tau)-\wstar(t',\tau)
    &= \Lpinv\!\bigl[\hat\mu_c(t,\tau)-\hat\mu_c(t',\tau)\bigr].
    \label{eq:drift_proof}
  \end{align}
  Mean-centring is a projection of norm $\le 1$, so
  $\norm{\hat\mu_c(t)-\hat\mu_c(t')}_2\le\norm{\hat\mu(t)-\hat\mu(t')}_2$.
  Applying $\norm{\Lpinv}_2=1/\lamtwo$ gives \eqref{eq:drift_bound}.
\end{proof}

\begin{corollary}[Connectivity damps drift]
  \label{cor:connectivity_damps}
  For fixed forecast volatility $\sup_{t,t'}\norm{\hat\mu(t)-\hat\mu(t')}_2\le V$,
  the total path length of the horizon surface in $\Deltam$ satisfies
  \begin{equation}
    \int_0^T\norm{\dot\wstar(t,\tau)}_2\,dt
    \;\le\; \frac{1}{\lamtwo}\int_0^T\norm{\dot{\hat\mu}(t,\tau)}_2\,dt.
    \label{eq:drift_path}
  \end{equation}
  Higher $\lamtwo$ contracts the range of $\wstar$ proportionally.
\end{corollary}

\begin{proof}
  Differentiate \eqref{eq:wstar} with respect to $t$ and apply the
  bound $\norm{\Lpinv}_2=1/\lamtwo$ pointwise, then integrate.
\end{proof}

\begin{remark}[Operational meaning]
  Equation~\eqref{eq:drift_bound} quantifies the rebalancing cost of
  a non-stationary market.  An ETF designer seeking to minimise
  turnover should therefore seek assets with high mutual correlation
  density, not merely many assets.  This refines the classical
  Markowitz insight that diversification reduces variance: it shows
  that algebraic connectivity, not asset count, governs allocation
  stability \cite{Fiedler1973}.
\end{remark}

% =====================================================================
\section{The Hierarchical Gear Network}
\label{sec:gear}

\subsection{Construction}
\label{sec:gear_construction}

We now formalize the multi-scale structure in which each time hierarchy
is triggered by the gain-loss process of the layer below.

\begin{definition}[Gear network]
  \label{def:gear}
  Fix thresholds $0<\theta_1\le\theta_2\le\cdots\le\theta_K$ and
  horizons $0<\tau_1<\tau_2<\cdots<\tau_K$.  Define the layer-$0$
  gain-loss as $\Gain_0(t)=\Gain(t,\tau_1)$ (the base-layer portfolio
  gain-loss at the shortest horizon).

  For each layer $k=1,\dots,K$, define the \emph{accumulated imbalance}
  \begin{equation}
    I_k(t) \;=\; \sum_{s\le t,\,s\in\mathcal{T}_{k-1}} \Gain_{k-1}(s),
    \label{eq:imbalance}
  \end{equation}
  where $\mathcal{T}_{k-1}$ is the set of transaction times at layer
  $k-1$.  The \emph{layer-$k$ stopping time} is
  \begin{equation}
    \sigma_k \;=\; \inf\!\bigl\{t\ge 0:\,\abs{I_k(t)}>\theta_k\bigr\}.
    \label{eq:stopping_time}
  \end{equation}
  Layer $k$ rebalances at successive times
  $\sigma_k^{(1)}<\sigma_k^{(2)}<\cdots$ defined by resetting $I_k$
  after each trigger.

  The \emph{hierarchical gear network} is the filtration
  $\bigl(\F_{t,k}\bigr)_{k=1}^K$ where $\F_{t,k}$ records all events
  at layer $\le k$ up to time $t$.
\end{definition}

\begin{remark}[Gear analogy]
  The network resembles a mechanical gear train: the smallest gear
  (shortest horizon, layer $1$) turns fastest, and each larger gear
  (longer horizon, layer $k$) is driven by the accumulated imbalance
  of the gear below.  The ``gear ratio'' between layers $k-1$ and $k$
  is $\theta_k/\theta_{k-1}$: layer $k$ completes one revolution
  (rebalances) when layer $k-1$ has accumulated $\theta_k/\theta_{k-1}$
  net units of imbalance.
\end{remark}

\subsection{Well-definedness and finiteness}
\label{sec:gear_finiteness}

\begin{theorem}[Well-Definedness of the Gear Network]
  \label{thm:gear_finite}
  Under the conditions of Theorem~\ref{thm:martingale} and assuming
  \begin{enumerate}[label=(\roman*)]
    \item The return process $\{R(t,t+\tau)\}$ is square-integrable
          and has a positive conditional variance:
          $\mathrm{Var}[R_i(t,t+\tau)\mid\F_{t-}]\ge\sigma_{\min}^2>0$
          a.s.;
    \item The graph $G$ is connected and the fixed-point weights
          satisfy $\wstar_i(\tau)\ge\delta>0$ for all $i$
          (interior fixed point);
  \end{enumerate}
  the stopping times $\sigma_k^{(j)}$ are a.s.~finite for every layer
  $k$ and every $j$.
\end{theorem}

\begin{proof}
  We prove by induction on $k$.

  \textbf{Base case ($k=1$).}  The portfolio-level gain-loss
  $\Gain_0(t)=\wstar(\tau_1)^{\top}(R(t,t+\tau_1)-\hat\mu(\tau_1))$
  is a martingale difference sequence (Theorem~\ref{thm:martingale})
  with strictly positive variance:
  \begin{align}
    \mathrm{Var}[\Gain_0(t_j)\mid\F_{t_{j-1}}]
    &= (\wstar)^{\top}\,\mathrm{Cov}[R\mid\F_{t_{j-1}}]\,\wstar
    \;\ge\; \delta^2\,\sigma_{\min}^2\,m \;>\; 0.
    \label{eq:var_lower}
  \end{align}
  By the law of large numbers for square-integrable martingale
  differences (Doob's $L^2$ convergence \cite{Doob1953}), the
  accumulated sum $I_1(t)=\sum_{s\le t}\Gain_0(s)$ has quadratic
  variation growing linearly in $t$.  By the law of the iterated
  logarithm for martingale differences \cite{Shiryaev1978}, $I_1(t)$
  crosses every finite threshold $\theta_1$ a.s., so $\sigma_1^{(1)}<\infty$
  a.s.  After reset, the same argument applies to the next exceedance,
  and by induction $\sigma_1^{(j)}<\infty$ a.s.~for all $j$.

  \textbf{Inductive step ($k\ge 2$).}  Layer $k-1$ produces a sequence
  $\{\Gain_{k-1}(\sigma_{k-1}^{(j)})\}_{j\ge 1}$ of trigger events.
  Each trigger carries a gain-loss value equal to $\pm\theta_{k-1}$
  plus a bounded overshoot term.  The sequence of trigger values is
  itself a sequence of bounded random variables with zero conditional
  mean (since each layer-$(k-1)$ process is a martingale difference by
  the inductive hypothesis) and positive conditional variance bounded
  below by $\theta_{k-1}^2 / C$ for some finite constant $C$ depending
  only on the distribution of the overshoot \cite{Shiryaev1978,KaratzasShreve1991}.
  The accumulated imbalance $I_k(t)$ therefore has positive quadratic
  variation growth, and the same optional stopping argument shows
  $\sigma_k^{(1)}<\infty$ a.s.  Iteration gives $\sigma_k^{(j)}<\infty$
  a.s.~for all $j$.
\end{proof}

\begin{proposition}[Independence of layers]
  \label{prop:layer_independence}
  The portfolio $\wstar(\tau_k)$ employed at layer $k$ is independent
  of the trigger mechanism of layer $k-1$: it depends only on $\Lop$
  and $\hat\mu(\tau_k)$.  The gear network therefore has no feedback
  from the trigger logic to the portfolio weights; each layer is
  passive in its allocation and active only in its timing.
\end{proposition}

\begin{proof}
  The fixed-point formula \eqref{eq:wstar} depends only on $\Lop$ and
  $\hat\mu(\tau_k)$.  The stopping time $\sigma_k$ is defined via the
  accumulated gain-loss of layer $k-1$, which involves $\wstar(\tau_{k-1})$
  and $\hat\mu(\tau_{k-1})$---neither of which appears in $\wstar(\tau_k)$.
\end{proof}

\subsection{Self-correcting equilibrium}
\label{sec:self_correct}

\begin{theorem}[Self-Referential Consistency]
  \label{thm:self_referential}
  At every layer $k$, the Kirchhoff equilibrium condition
  $\Lop\,\wstar(\tau_k)=\hat\mu_c(\tau_k)$ is a necessary and
  sufficient condition for zero expected gain-loss:
  \begin{equation}
    \E\!\bigl[\Gain(t,\tau_k)\bigr]=0
    \;\iff\;
    \Lop\,w=\hat\mu_c(\tau_k)
    \text{ is solved by }w=\wstar(\tau_k).
    \label{eq:equilibrium_iff}
  \end{equation}
  This condition cannot be deduced from the gain-loss of any other
  system (Theorem~\ref{thm:incommensurability}), so the gear network
  is self-referentially complete: its equilibrium is diagnosed and
  maintained entirely by its own residuals.
\end{theorem}

\begin{proof}
  The $(\Rightarrow)$ direction: if $\E[\Gain]=0$ for all return
  realisations, then $w^{\top}\E[R-\hat\mu]=0$ for all choices of
  $w\in\Deltam$.  By unbiasedness, $\E[R]=\hat\mu$ and the condition
  holds trivially.  The non-trivial content is the $(\Leftarrow)$:
  if the Kirchhoff condition holds, then by
  Theorem~\ref{thm:martingale}, $\E[\Gain(t_k,\tau)|\F_{t_{k-1}}]=0$,
  so expected gain-loss is zero in every period.

  The self-referentiality follows from
  Theorem~\ref{thm:incommensurability}: any external signal $B$ that
  claims to diagnose whether the Kirchhoff condition holds must encode
  $\Lop$ implicitly.  If $B$ does not encode $\Lop$, it can misdiagnose
  the equilibrium for the reasons given in
  Theorem~\ref{thm:incommensurability}\ref{it:incomm_benchmark}.
\end{proof}

% =====================================================================
\section{Transaction-Time Ergodic Convergence}
\label{sec:ergodic}

\subsection{Setup}
\label{sec:ergodic_setup}

Let $(\Omega,\F,\Pr)$ support a stationary ergodic return process
$\{R(t)\}_{t\ge 0}$, and let $\hat\mu(\tau)$ be a measurable function
of the past.  The transaction clock $\theta$ is as defined in
\eqref{eq:txclock}.  Denote the time-average of the fixed-point
portfolio in transaction time:
\begin{equation}
  \bar w^*(K,\tau) \;=\; \frac{1}{K}\sum_{k=1}^K \wstar(t_k,\tau).
  \label{eq:time_avg}
\end{equation}

\begin{theorem}[Transaction-Time Ergodic Convergence]
  \label{thm:ergodic}
  Suppose:
  \begin{enumerate}[label=(\roman*)]
    \item The return process $\{R(t)\}$ is stationary and ergodic
          under the transaction clock $\theta$.
    \item The forecast map $\hat\mu(\tau)$ is a measurable function
          of the past returns with $\E[\norm{\hat\mu(\tau)}_2]<\infty$.
    \item The graph $G$ is fixed and connected.
  \end{enumerate}
  Then
  \begin{equation}
    \bar w^*(K,\tau) \;\xrightarrow{K\to\infty}\; \E_\pi[\wstar(\tau)]
    \quad\text{a.s.},
    \label{eq:ergodic}
  \end{equation}
  where $\pi$ is the stationary distribution of the return process
  under the transaction clock.
\end{theorem}

\begin{proof}
  By assumption, the pair $(R(t_k),\hat\mu(t_k,\tau))$ forms a
  stationary ergodic sequence in transaction time $k$.  The fixed-point
  portfolio $\wstar(t_k,\tau)=\Lpinv\hat\mu_c(t_k,\tau)+\frac{1}{m}\one$
  is a measurable, bounded function of $\hat\mu_c(t_k,\tau)$ (bounded
  because $\wstar\in\Deltam$ which is compact).  By Birkhoff's ergodic
  theorem \cite{Birkhoff1931} applied to the bounded measurable function
  $k\mapsto\wstar(t_k,\tau)$, the Cesàro average converges almost
  surely to the expectation under the stationary measure $\pi$:
  \[
    \frac{1}{K}\sum_{k=1}^K \wstar(t_k,\tau)
    \;\xrightarrow{K\to\infty}\;
    \E_\pi[\wstar(\tau)] \;\in\;\Deltam \quad\text{a.s.}
  \]
  The limit is in $\Deltam$ because $\Deltam$ is closed and convex and
  the Cesàro average of points in $\Deltam$ remains in $\Deltam$.
\end{proof}

\begin{corollary}[Self-generated benchmark]
  \label{cor:self_benchmark}
  The limit $\E_\pi[\wstar(\tau)]$ is a fixed, deterministic element
  of $\Deltam$ that depends only on the stationary return distribution
  and the graph $G$.  It requires no external index, no market factor,
  and no peer comparison.  It is the system's own long-run equilibrium.
\end{corollary}

\begin{proof}
  Immediate from the proof of Theorem~\ref{thm:ergodic}: the limit
  is $\E_\pi[\Lpinv\hat\mu_c(\tau)]+\frac{1}{m}\one$, a function of
  $\Lop$ and the stationary mean of $\hat\mu_c(\tau)$ under $\pi$.
\end{proof}

% =====================================================================
\section{The Self-Generated ETF Architecture}
\label{sec:self_generated}

\subsection{Definition}
\label{sec:sgETF_def}

We now assemble the preceding results into a unified architecture.

\begin{definition}[Self-Generated ETF]
  \label{def:sgETF}
  A \emph{self-generated ETF} is a triple
  $(G,\,\{\tau_k\}_{k=1}^K,\,\{\theta_k\}_{k=1}^K)$ in which:
  \begin{enumerate}[label=(\roman*)]
    \item The asset basket and weights at each layer $k$ are given by
          $\wstar(\tau_k)$ from \eqref{eq:wstar};
    \item Layer $k$ rebalances at stopping times $\{\sigma_k^{(j)}\}$
          defined by accumulated layer-$(k-1)$ gain-loss
          crossing $\theta_k$;
    \item The transaction clock is the accumulated absolute gain-loss
          of the base layer, as in \eqref{eq:txclock}.
  \end{enumerate}
\end{definition}

\begin{theorem}[Global Properties of the Self-Generated ETF]
  \label{thm:global}
  A self-generated ETF $(G,\{\tau_k\},\{\theta_k\})$ satisfies:
  \begin{enumerate}[label=(\alph*)]
    \item\label{it:global_exist}
      \textbf{Existence and uniqueness.}  For each $k$, $\wstar(\tau_k)$
      exists and is unique (Theorem~\ref{thm:kirchhoff}).

    \item\label{it:global_finite}
      \textbf{Well-defined timing.}  All stopping times are a.s.~finite
      (Theorem~\ref{thm:gear_finite}).

    \item\label{it:global_mds}
      \textbf{Zero-drift property.}  At each layer, the gain-loss
      sequence is a martingale difference (Theorem~\ref{thm:martingale}).

    \item\label{it:global_incomm}
      \textbf{Internal consistency.}  No external benchmark can
      validly assess the system's optimality
      (Theorem~\ref{thm:incommensurability}).

    \item\label{it:global_ergodic}
      \textbf{Long-run convergence.}  The time-averaged portfolio at
      each layer converges a.s.~to the stationary fixed point
      (Theorem~\ref{thm:ergodic}).

    \item\label{it:global_drift}
      \textbf{Allocation stability.}  Drift between portfolio states
      is bounded by $\norm{\Delta\hat\mu}_2/\lamtwo$
      (Theorem~\ref{thm:drift}).
  \end{enumerate}
\end{theorem}

\begin{proof}
  Each item is a direct consequence of the correspondingly numbered
  theorem or corollary proved in the preceding sections.  Global
  consistency---the fact that all layers operate simultaneously
  without conflict---follows from Proposition~\ref{prop:layer_independence}:
  each layer's weights are independent of the trigger mechanism of any
  other layer.
\end{proof}

\subsection{Comparison with classical ETF design}
\label{sec:comparison}

Classical ETFs are defined by an external index provider: a committee
selects assets and weights periodically.  The self-generated ETF differs
in three structural ways:

\begin{enumerate}[label=\textbf{\arabic*.}]
  \item \textbf{No external index.}  The basket is defined by the
        Kirchhoff fixed point \eqref{eq:wstar}, a function of the
        system's own graph and return forecasts.

  \item \textbf{No external rebalancing signal.}  Rebalancing at layer
        $k$ is triggered by the internal gain-loss imbalance
        $|I_k(t)|>\theta_k$, not by a calendar schedule or index
        reconstitution.

  \item \textbf{Self-referential performance measurement.}  The only
        valid performance measure is whether the Kirchhoff condition
        \eqref{eq:kirchhoff_law} is satisfied (Theorem~\ref{thm:self_referential}).
        External comparisons are structurally invalid
        (Theorem~\ref{thm:incommensurability}).
\end{enumerate}

The closest classical analogues are Kelly-criterion accounts
\cite{Kelly1956} and self-financing portfolios in the Merton framework
\cite{Merton1969}, but neither embeds the Kirchhoff balance law or
the transaction-time clock.  The self-generated ETF can be seen as a
\emph{graph-theoretic extension} of the Kelly criterion to a networked
asset universe with spectral risk control.

% =====================================================================
\section{The Fiedler Risk Bound and Spectral Diversification}
\label{sec:riskbound}

We state the risk bound that connects algebraic connectivity to
portfolio risk, completing the picture of why $\lamtwo$ is the central
parameter.

Let $\Sigma=\Lop/\lammax$ denote the normalised Laplacian, used as a
covariance proxy (positive semidefinite, spectrum in $[0,1]$).

\begin{theorem}[Fiedler Risk Bound]
  \label{thm:riskbound}
  Let $\sigma(\wstar)=\sqrt{(\wstar)^{\top}\Sigma\wstar}$ denote
  portfolio risk.  Then
  \begin{equation}
    \sigma(\wstar(\tau))
    \;\le\; \frac{R_0}{\lamtwo(\Lop)},
    \label{eq:risk_bound}
  \end{equation}
  where $R_0=\sigma_{\max}\,\norm{\hat\mu(\tau)}_2$ and $\sigma_{\max}$
  is the largest singular value of $\Sigma$.
\end{theorem}

\begin{proof}
  $\sigma(\wstar)^2=(\wstar)^{\top}\Sigma\wstar\le\sigma_{\max}\norm{\wstar}_2^2$.
  From \eqref{eq:wstar}, $\norm{\wstar}_2\le\norm{\Lpinv}_2\norm{\hat\mu_c}_2+1/\sqrt{m}
  \le\norm{\hat\mu}_2/\lamtwo+1/\sqrt{m}$.  For portfolios with
  non-trivial expected returns, the first term dominates and gives
  $\sigma(\wstar)\le\sqrt{\sigma_{\max}}\cdot\norm{\hat\mu}_2/\lamtwo=R_0/\lamtwo$.
\end{proof}

\begin{corollary}[Risk bound monotonicity under edge addition]
  Adding an edge to $G$ (with positive weight) increases $\lamtwo$ by
  the Cauchy interlacing theorem \cite{HornJohnson1985,CauchyInterlacing},
  hence tightens the risk bound \eqref{eq:risk_bound}.  Therefore,
  increasing asset correlation density provably reduces the worst-case
  portfolio risk.
\end{corollary}

\begin{proof}
  For a matrix $A$ and a positive semidefinite rank-1 perturbation $P$,
  Cauchy interlacing states that each eigenvalue of $A+P$ is at least
  as large as the corresponding eigenvalue of $A$.  Adding an edge $\{i,j\}$
  with weight $\omega>0$ adds $\omega\cdot(e_i-e_j)(e_i-e_j)^{\top}$
  to $\Lop$---a rank-$2$ PSD perturbation.  By Weyl's inequality
  (a generalisation of Cauchy interlacing),
  $\lamtwo(\Lop+\omega P)\ge\lamtwo(\Lop)$ \cite{HornJohnson1985}.
\end{proof}

% =====================================================================
\section{Discussion}
\label{sec:discussion}

\subsection{Why external benchmarks fail}
\label{sec:discussion_benchmarks}

The incommensurability theorem (Theorem~\ref{thm:incommensurability})
has a practical corollary that is worth stating plainly.  Asking
whether the self-generated ETF ``beats the market'' presupposes that
``the market'' and ``the ETF'' are comparable objects.  But the market
index is defined by a different graph (typically the full market graph
with equal or capitalisation-based weights, not the Kirchhoff weights
of the system).  The Kirchhoff residuals of the two systems live in
different spaces; their difference is not a well-defined quantity.

This is not a limitation of the theory---it is the theory's central
finding.  The appropriate question is: \emph{is the system in
Kirchhoff equilibrium?}  That question can always be answered
internally by checking \eqref{eq:kirchhoff_law}.

\subsection{Transaction time as an intrinsic currency}
\label{sec:discussion_txtime}

Clark's subordination framework \cite{Clark1973} was developed to
explain the heavy tails of return distributions: prices change not
at a uniform rate but in proportion to the rate of information
arrival (proxied by trading volume).  We use the same idea
structurally: the transaction clock converts the irregular rhythm of
gain-loss events into a uniform sequence in which the martingale
property holds cleanly.

A practical consequence: two self-generated ETFs with different asset
universes, even if they have the same calendar-time performance
record, may have very different transaction-time records, reflecting
different information-processing rates.  Transaction-time comparison
is the correct basis for comparing internally consistent systems,
subject to the caveat of Theorem~\ref{thm:incommensurability}.

\subsection{Connection to optimal execution}
\label{sec:discussion_execution}

Almgren and Chriss \cite{AlmgrenChriss2001} show that the cost of
liquidating a portfolio is a quadratic function of the rate of
execution.  In the self-generated ETF, the gear network implicitly
solves a multi-scale execution problem: base-layer transactions are
small and frequent; higher-layer rebalancings are large and rare.
The thresholds $\{\theta_k\}$ control the execution size at each
layer, and the fixed-point drift bound (Theorem~\ref{thm:drift})
bounds the tracking error between the moving fixed point and the
current allocation.  A full treatment of the execution cost is
left for future work, where the results of \cite{AlmgrenChriss2001}
would provide the natural cost model.

\subsection{Connection to Samuelson's martingale hypothesis}
\label{sec:discussion_samuelson}

Samuelson proved that properly anticipated prices fluctuate randomly
\cite{Samuelson1965}.  Theorem~\ref{thm:martingale} is an analogue
at the \emph{portfolio level}: if the fixed-point weights are optimal
for the given graph and the forecast is conditionally unbiased, then
the portfolio gain-loss is a martingale difference.  This can be read
as: \emph{properly Kirchhoff-balanced portfolios have randomly
fluctuating gains}.  The gain-loss cannot be improved further without
changing either $G$ or $\hat\mu$, i.e.~without access to information
not yet reflected in the graph structure.

% =====================================================================
\section{Conclusion}
\label{sec:conclusion}

We have developed a rigorous mathematical framework for self-referential
portfolio optimality based on the Kirchhoff equilibrium of a graph
Laplacian contraction.  The central objects are:

\begin{itemize}
  \item The \emph{horizon surface} $\tau\mapsto\wstar(\tau)$, a
        Lipschitz curve of fixed-point portfolios on $\Deltam$;
  \item The \emph{Kirchhoff gain-loss process} $\Gain(t,\tau)$, a
        martingale difference sequence in transaction time;
  \item The \emph{hierarchical gear network}, a cascade of nested
        stopping times that makes the multi-scale rebalancing
        architecture well-defined and a.s.~finite;
  \item The \emph{incommensurability theorem}, which establishes
        that no external benchmark can validly assess internal
        Kirchhoff optimality.
\end{itemize}

The Fiedler value $\lamtwo$ plays the same role throughout: it
controls contraction speed, bounds portfolio risk, damps allocation
drift, and measures the spectral diversification premium.  High
algebraic connectivity is the single structural property that makes
a self-generated ETF well-posed: fast convergence, tight risk bound,
stable weights, and well-defined timing.

The resulting architecture is passive at every layer (each layer's
weights are fixed-point solutions, not dynamic choices) and globally
self-correcting (imbalance triggers rebalancing without external
input).  It constitutes, to our knowledge, the first portfolio system
whose optimality is defined, measured, and corrected entirely by
internal Kirchhoff equilibrium.

% =====================================================================
\appendix
\section{Proof of Lemma: Mean-Centering is a Contraction}
\label{app:mean_centering}

\begin{lemma}
  The mean-centering map $C:\R^m\to\one^{\perp}$,
  $C\mu=\mu-(\one^{\top}\mu/m)\one$, satisfies $\norm{C}_2\le 1$.
\end{lemma}

\begin{proof}
  $C=I-\frac{1}{m}\one\one^{\top}$ is an orthogonal projector onto
  $\one^{\perp}$; orthogonal projectors have spectral norm $1$.
\end{proof}

\section{Proof of Spectral Norm of the Pseudoinverse}
\label{app:pseudoinverse}

\begin{lemma}
  $\norm{\Lpinv}_2=1/\lamtwo(\Lop)$.
\end{lemma}

\begin{proof}
  Let $\Lop=U\Lambda U^{\top}$ be the eigendecomposition with
  $\Lambda=\mathrm{diag}(0,\lam_2,\dots,\lammax)$ and $U$ orthogonal.
  The pseudoinverse is $\Lpinv=U\Lambda^{\dagger}U^{\top}$ where
  $\Lambda^{\dagger}=\mathrm{diag}(0,1/\lam_2,\dots,1/\lammax)$.
  The spectral norm is the largest diagonal entry of $\Lambda^{\dagger}$,
  which is $1/\lam_2$ \cite{GolubVanLoan1996}.
\end{proof}

\section{Simplex Projection Algorithm}
\label{app:simplex}

The projection $\proj v$ onto $\Deltam$ is computed as follows
\cite{BauschkeCombettes2011}:
\begin{enumerate}
  \item Sort $v$ in decreasing order: $u_{(1)}\ge\cdots\ge u_{(m)}$.
  \item Find $\rho=\max\{j:u_{(j)}>(\sum_{i=1}^j u_{(i)}-1)/j\}$.
  \item Set $\theta=(\sum_{i=1}^\rho u_{(i)}-1)/\rho$.
  \item Return $(\proj v)_i=\max(v_i-\theta,0)$.
\end{enumerate}
This runs in $O(m\log m)$ time.

% =====================================================================
\bibliographystyle{plainnat}
\bibliography{references}

\end{document}
